\documentclass[14pt,a4paper,oneside]{extbook}
\usepackage[utf8]{vietnam}
\usepackage{amsmath,amssymb,mathrsfs,fancyhdr,enumerate,multirow,makecell}
\usepackage{tkz-euclide,tikz-3dplot,tikz,tkz-tab}
\usetikzlibrary{arrows.meta,calc,intersections,angles,quotes,backgrounds,shapes.geometric,patterns,through,shadows,positioning}
\usepackage{yhmath}
\usepackage[top=0.7cm, bottom=0.7cm, left=1.0cm, right=1.0cm, headheight=16pt, headsep=3mm, footskip=5mm]{geometry}
\usepackage[hidelinks,unicode]{hyperref}
\usepackage{qrcode}
\usepackage[many]{tcolorbox}
\usepackage{fontawesome}

\pagestyle{fancy}
\fancyhf{}
\fancyhead[L]{\footnotesize\bfseries\color{teal!85!black}\faBook~SỔ TAY KIẾN THỨC TOÁN 5}
\fancyhead[R]{\footnotesize\bfseries\color{blue!85!black}Lớp Toán Thầy Lĩnh -- loptoanthaylinh.vn}
\fancyfoot[C]{\footnotesize\bfseries\color{gray!80!black}-- Sổ tay kiến thức: Bài 60. Quãng đường, thời gian của một chuyển động đều --}
\renewcommand{\headrulewidth}{0.4pt}

\begin{document}

% --- BANNER TIÊU ĐỀ BÀI HỌC ---
\noindent
\begin{tcolorbox}[
  enhanced,
  colback=teal!8!white,
  colframe=teal!80!black,
  arc=3mm,
  boxrule=1.1pt,
  width=\linewidth,
  halign=center,
  top=2.2mm, bottom=2.2mm,
  left=3mm, right=3mm,
  drop fuzzy shadow
]
  {\large\bfseries\color{teal!95!black}\faBookmark\ BÀI 60. QUÃNG ĐƯỜNG, THỜI GIAN CỦA CHUYỂN ĐỘNG ĐỀU}\\[2pt]
  {\scriptsize\color{black!75}\itshape Công thức s = v x t và t = s : v, xử lý thời điểm khởi hành, thời điểm đến trên 1 trang A4}
\end{tcolorbox}

\vspace{3.2mm}
% --- KHỐI 1: CÔNG THỨC QUÃNG ĐƯỜNG VÀ THỜI GIAN ---
\noindent
\begin{minipage}[t]{0.50\textwidth}
  \begin{tcolorbox}[
    enhanced,
    title={\footnotesize\bfseries\color{white}\faRoad\ 1. Công thức Quãng đường ($s$)},
    colback=blue!2!white,
    colframe=blue!75!black,
    arc=2.5mm, boxrule=0.8pt, top=1.8mm, bottom=1.8mm, left=2.2mm, right=2.2mm
  ]
    \footnotesize
    \begin{tcolorbox}[colback=white, colframe=blue!75!black, boxrule=0.5pt, arc=1.5mm, halign=center, top=1mm, bottom=1mm]
      {\scriptsize\bfseries\color{blue!90!black}$s = v \times t$}
    \end{tcolorbox}
    \textit{Muốn tính quãng đường, ta lấy vận tốc nhân với thời gian.}\\
    $\bullet$ Đơn vị phải tương ứng: Nếu $v$ là $\text{km/h}$, $t$ là $\text{giờ}$ thì $s$ là $\text{km}$.\\
    \textit{VD:} Xe chạy $45\text{ km/h}$ trong $3\text{ giờ}$:\\
    $s = 45 \times 3 = \mathbf{135\text{ km}}$.
  \end{tcolorbox}
\end{minipage}%
\hfill
\begin{minipage}[t]{0.48\textwidth}
  \begin{tcolorbox}[
    enhanced,
    title={\footnotesize\bfseries\color{white}\faHourglassHalf\ 2. Công thức Thời gian ($t$)},
    colback=teal!2!white,
    colframe=teal!75!black,
    arc=2.5mm, boxrule=0.8pt, top=1.8mm, bottom=1.8mm, left=2mm, right=2mm
  ]
    \footnotesize
    \begin{tcolorbox}[colback=white, colframe=teal!75!black, boxrule=0.5pt, arc=1.5mm, halign=center, top=1mm, bottom=1mm]
      {\scriptsize\bfseries\color{teal!95!black}$t = s : v$}
    \end{tcolorbox}
    \textit{Muốn tính thời gian, ta lấy quãng đường chia cho vận tốc.}\\
    $\bullet$ Kết quả dạng số thập phân cần \textbf{đổi ra giờ và phút}.\\
    \textit{VD:} $t = 120 : 50 = 2,4\text{ giờ}$.\\
    Đổi: $2,4\text{h} = 2\text{h} + 0,4 \times 60\text{p} = \mathbf{2\text{h } 24\text{p}}$.
  \end{tcolorbox}
\end{minipage}

\vspace{3mm}

% --- KHỐI 2: DẠNG TOÁN KHỞI HÀNH VÀ ĐẾN NƠI ---
\noindent
\begin{tcolorbox}[
  enhanced,
  title={\footnotesize\bfseries\color{white}\faCalendarCheckO\ 3. Mối liên hệ Thời điểm đến và Thời điểm đi},
  colback=purple!2!white,
  colframe=purple!70!black,
  arc=2.5mm, boxrule=0.8pt, top=1.8mm, bottom=1.8mm, left=2.5mm, right=2.5mm
]
  \footnotesize
  Một người đi xe máy từ A lúc $7\text{ giờ } 15\text{ phút}$, đến B với vận tốc $40\text{ km/h}$. Quãng đường AB dài $100\text{ km}$. Dọc đường nghỉ $15\text{ phút}$. Hỏi người đó đến B lúc mấy giờ?\\
  $\triangleright$ Thời gian thực tế đi trên đường: $t = 100 : 40 = 2,5\text{ giờ} = 2\text{ giờ } 30\text{ phút}$.\\
  $\triangleright$ Thời điểm đến B là: $7\text{h } 15\text{p} + 2\text{h } 30\text{p} + 15\text{p (nghỉ)} = \mathbf{10\text{ giờ } 00\text{ phút}}$.
\end{tcolorbox}

\vspace{3mm}

% --- KHỐI 3: SƠ ĐỒ TƯ DUY ---
\noindent
\begin{tcolorbox}[
  enhanced,
  title={\footnotesize\bfseries\color{white}\faSitemap\ Sơ đồ tư duy: Ba đại lượng chuyển động đều},
  colback=yellow!2!white,
  colframe=teal!85!black,
  arc=3mm, boxrule=0.9pt, top=1.8mm, bottom=1.8mm, left=1.5mm, right=1.5mm,
  halign=center
]
  \begin{tikzpicture}[
    scale=0.92,
    every node/.style={transform shape},
    b1/.style={rectangle, rounded corners=4pt, fill=blue!15, draw=blue!70!black, text=blue!90!black, font=\bfseries\scriptsize, inner sep=4pt, align=center},
    b2/.style={rectangle, rounded corners=4pt, fill=teal!15, draw=teal!70!black, text=teal!90!black, font=\bfseries\scriptsize, inner sep=4pt, align=center},
    b3/.style={rectangle, rounded corners=4pt, fill=green!15, draw=green!70!black, text=green!90!black, font=\bfseries\scriptsize, inner sep=4pt, align=center}
  ]
    \node[b1] (n1) at (-5.8, 0) {Tìm $s$: $s = v \times t$\\Đổi $t$ ra giờ};
    \node[b2] (n2) at (0, 0) {Tìm $v$: $v = s : t$\\Đơn vị $\text{km/h}$};
    \node[b3] (n3) at (5.8, 0) {Tìm $t$: $t = s : v$\\Đổi $t$ ra giờ và phút};

    \draw[teal!85!black, thick, <->] (n1) -- (n2);
    \draw[teal!85!black, thick, <->] (n2) -- (n3);
  \end{tikzpicture}
\end{tcolorbox}

\vspace{3mm}

% --- KHỐI 4: ĐỔI PHẦN THẬP PHÂN RA PHÚT & BẪY TRẮC NGHIỆM ---
\noindent
\begin{minipage}[t]{0.54\textwidth}
  \begin{tcolorbox}[
    enhanced,
    title={\footnotesize\bfseries\color{white}\faBolt\ Cách đổi phần thập phân của giờ ra phút},
    colback=green!3!white,
    colframe=green!60!black,
    arc=2.5mm, boxrule=0.8pt, top=1.8mm, bottom=1.8mm, left=2.2mm, right=2.2mm
  ]
    \footnotesize
    Lấy phần thập phân \textbf{nhân với 60}:\\
    $\bullet$ $1,2\text{ giờ} = 1\text{ giờ} + (0,2 \times 60)\text{ phút} = 1\text{ giờ } 12\text{ phút}$.\\
    $\bullet$ $2,75\text{ giờ} = 2\text{ giờ} + (0,75 \times 60)\text{ phút} = 2\text{ giờ } 45\text{ phút}$.\\
    $\bullet$ $3,3\text{ giờ} = 3\text{ giờ } 18\text{ phút}$ (vì $0,3 \times 60 = 18$).
  \end{tcolorbox}
\end{minipage}%
\hfill
\begin{minipage}[t]{0.44\textwidth}
  \begin{tcolorbox}[
    enhanced,
    title={\footnotesize\bfseries\color{white}\faExclamationTriangle\ Mẹo tránh bẫy trắc nghiệm},
    colback=red!2!white,
    colframe=red!70!black,
    arc=2.5mm, boxrule=0.8pt, top=1.8mm, bottom=1.8mm, left=2.2mm, right=2.2mm
  ]
    \footnotesize
    $\bullet$ Rất nhiều học sinh đổi: $2,4\text{ giờ} = 2\text{ giờ } 40\text{ phút}$ hoặc $2\text{ giờ } 4\text{ phút}$ (\textbf{SAI}). Phải lấy $0,4 \times 60 = \mathbf{24\text{ phút}}$.\\
    $\bullet$ Quên cộng thời gian nghỉ khi tính thời điểm đến nơi.
  \end{tcolorbox}
\end{minipage}

\end{document}
